% ATM'22 dataset -- pulled in by sections/03-dataset.tex

\subsection{The ATM'22 dataset}  % first dataset
\label{sec:atm22}                % lets other sections write "see Section~\ref{sec:atm22}"

The Multi-site, Multi-domain Airway Tree Modeling dataset (ATM'22) was released
as part of an official challenge held in conjunction with MICCAI 2022
\cite{zhang2023atm22}. It targets the binary segmentation of the pulmonary
airway tree from chest CT, and at the time of its release it was the largest
publicly available airway dataset with complete voxel-wise annotations. The
reference benchmark that preceded it, EXACT'09, provided only 40 scans and
never released its manual annotations, which had long limited the training and
fair evaluation of data-driven methods.

The dataset contains 500 chest CT scans, partitioned into 300 training, 50
validation and 150 test cases. The test set is held privately by the organizers
and is accessible only through the submission of a Docker container, which
guarantees the fairness and reproducibility of the reported results. The scans
originate from three sources, as summarized in Table~\ref{tab:atm22_split}.

\begin{table}[htbp]              % htbp = let LaTeX choose a good spot for the table
  \centering                     % centre the table on the page
  \caption{Composition of the ATM'22 dataset by split and by source site.}
  \label{tab:atm22_split}        % target of the \ref above
  \begin{tabular}{lcccc}         % 1 left-aligned column + 4 centred ones
    \toprule
    Split & Shanghai Chest Hospital & LIDC-IDRI & EXACT'09 & Total \\
    \midrule
    Training   & 226                  & 54 & 20 & 300 \\
    Validation & 20                   & 30 & -- & 50  \\
    Test       & 84 (58 COVID-19)     & 66 & -- & 150 \\
    \midrule
    Total      & 330                  & 150 & 20 & 500 \\
    \bottomrule
  \end{tabular}
\end{table}

For the 20 cases taken from EXACT'09, only the airway annotations are
distributed; the corresponding CT volumes must be obtained separately from the
original source.

The acquisition parameters are heterogeneous, which is the basis of the
multi-site character of the dataset. The scans were acquired on Philips iCT 256,
GE LightSpeed16 and Toshiba Aquilion scanners, the last of which appears only in
the test set. Each slice is $512 \times 512$ pixels, and the number of slices per
scan ranges from 157 to 1125, with a slice thickness between 0.45 and 1.00~mm
and an in-plane resolution between 0.500 and 0.919~mm. The health status of the
subjects ranges from healthy volunteers to patients with severe pulmonary
disease. The multi-domain character comes from the 58 COVID-19 scans included in
the test set: ground-glass opacity, consolidation and septal thickening alter
the appearance of the lung parenchyma and blur the airway lumen, so these cases
form a noisy domain on which the generalization ability of a model trained on
clean scans can be assessed.

The annotations were produced in two stages. A preliminary segmentation was
first obtained by ensembling, through majority voting, several airway
segmentation models trained on the BAS dataset. These preliminary masks were
then reviewed by three radiologists with more than five years of experience: two
of them independently refined and cross-checked each scan, and the most senior
radiologist validated the final annotation. The manual refinement required 60 to
90 minutes per scan and focused mainly on repairing breakages in the peripheral
bronchi and removing small leakages into the surrounding parenchyma, using the
original CT intensities and anatomical prior knowledge. The whole collection and
annotation process took the organizers approximately one year.

The challenge evaluates submissions along two complementary axes. Topological
completeness is measured by the tree length detected rate (TD) and the branch
detected rate (BD),
\begin{equation}                 % numbered equation, referable with \eqref{eq:td_bd}
  \mathrm{TD} = \frac{T_{det}}{T_{ref}},
  \qquad
  \mathrm{BD} = \frac{B_{det}}{B_{ref}},
  \label{eq:td_bd}
\end{equation}
where $T_{det}$ and $B_{det}$ denote respectively the total length and the
number of branches correctly detected in the prediction, and $T_{ref}$ and
$B_{ref}$ the corresponding quantities in the ground truth. Both metrics are
computed on the largest connected component of the prediction, a branch being
counted as detected when more than 80\% of its centerline voxels lie within the
ground truth. Topological correctness is measured by the Dice similarity
coefficient (DSC) and the precision, with sensitivity and specificity reported
as complementary indicators. The official ranking is based on the unweighted
mean of TD, BD, DSC and precision. This dual evaluation reflects the fact that
overlap-based measures alone are insufficient for tubular structures: a
prediction may achieve a high Dice score while missing a substantial number of
distal branches, since the peripheral bronchi account for only a small fraction
of the total airway volume.

